% ============================================================
% chapter_04_results.tex - Chapter 4, Results
%
% Purpose: reports the empirical performance of the LSEQM+DL framework
%          against 172 independent BMKG stations across the four
%          verification pillars + temporal-skill diagnostics.
% Sections: 4.1 Agreement Among the Three Datasets,
%           4.2 Stage-wise Headline Skill, 4.3 Stage-wise Attribution,
%           4.4 Threshold-Dependent Detection,
%           4.5 Pearson Correlation Across Stages,
%           4.6 Spatial Skill Distribution, 4.7 Per-Station Taylor Diagram
%           Analysis, 4.8 Temporal Stability Across Dekads,
%           4.9 Sensitivity Analyses, 4.10 Reproducibility Evaluation,
%           4.11 Chapter summary.
% Included by main_thesis.tex.
% License: CC BY 4.0 (see ../LICENSE).
% ============================================================

\chapter{RESULTS}\label{chap:results}

This chapter reports the empirical performance of the LSEQM+DL
framework (Linear Scaling, Empirical Quantile Mapping, Generalized
Pareto Distribution, and a Convolutional Neural Network refinement)
against 172 independent BMKG stations, pooled across
the 36 dekadal windows of 2001 to 2021. The results are organised
around the four-pillar verification framework introduced in
Section~\ref{sec:methods-metrics}, with the temporal-skill
diagnostic of Pearson correlation reported separately because of
the structural distinction that
Chapter~\ref{chap:discussion} develops.

Section~\ref{sec:results-intercomparison} first compares the three
datasets with each other, before any correction, and separates
their genuine disagreement from the calendar label each archive
assigns.
Section~\ref{sec:results-headline} presents the headline
stage-wise skill across the four pillars.
Section~\ref{sec:results-stages} attributes each metric movement
to the pipeline stage that produced it.
Section~\ref{sec:results-threshold} reports the
threshold-dependent detection results that the aggregate
event-detection numbers in
Section~\ref{sec:results-headline} mask.

Section~\ref{sec:results-correlation} reports the Pearson
correlation across stages and the paradox at the centre of the
discussion.
Section~\ref{sec:results-window} returns to the calendar theme of
Section~\ref{sec:results-intercomparison} at half-hourly
resolution, locating the optimal aggregation window; the two are
independent measurements of the same effect and agree. The
remaining sections cover the spatial skill
distribution (Section~\ref{sec:results-cqi}), per-station Taylor
analysis (Section~\ref{sec:results-taylor}), temporal stability
across dekads (Section~\ref{sec:results-temporal}), sensitivity
analyses on the framework parameters
(Section~\ref{sec:results-sensitivity}), and the reproducibility
evaluation that demonstrates end-to-end re-runnability
(Section~\ref{sec:results-reproducibility}).


% ============================================================
\section{Agreement Among the Three Datasets}\label{sec:results-intercomparison}
% ============================================================

The correction pairs a satellite product with a gridded gauge
reference and is judged against an independent station network, so
the three datasets are compared with each other before the
correction results are presented. Two questions are settled here:
how closely the three agree at the daily scale, and how much of any
disagreement belongs to the data rather than to the calendar label
each archive assigns. The corrected stages are then put through the
same comparison, which separates the two effects.

Every satellite-versus-BMKG correlation in this section is
constructed the same way, and differently from the rest of the
chapter. It uses daily paired values, taken as a single pooled
correlation over all station-days rather than sliced into
dekad-of-year windows, over the continuous GPM-era record from 2015
to 2021, against the $177$ BMKG stations that carry at least $100$
valid paired days in that window. The grid-to-grid
pairs reported below, IMERG-L against CPC-UNI and the corrected
fields against IMERG-F, are necessarily constructed differently
and are qualified where they are quoted. That count sits
between the $172$ stations used for the dekadal correction metrics
(Section~\ref{sec:results-headline}) and the $178$ used for the
half-hourly window diagnostic (Section~\ref{sec:results-window});
the three differ only in the coverage each analysis requires.
Full-record values are reported alongside throughout, because the
two eras differ substantially.

Because these values are constructed differently from the
correlations reported elsewhere in the thesis, they do not match
them, and the differences are systematic rather than
contradictory. Three constructions recur. The in-sample figure of
$0.348$ quoted in Section~\ref{sec:results-correlation} is measured
per pixel against the CPC-UNI calibration target, not against the
stations. Against the BMKG stations over the full record the same
corrected product gives $0.280$; slicing that record into the $36$
dekad-of-year windows used for the correction metrics removes the
seasonal cycle and lowers it to $0.235$; restricting instead to the
GPM era lowers it to $0.174$. Each step changes the construction
in more than one respect, stratification and pooling together, so
the values are not directly comparable; the reference, the record,
the stratification, the pooling, and the window label must all be
carried with any correlation quoted from this work.

\subsection{Agreement Under the Archived Date Labels}

Under the date label each archive carries, the three daily
correlations split sharply. The two UTC-dated products, IMERG-L and
CPC-UNI, agree with each other at $r = 0.57$, measured between the
two gridded archives at the $178$ station locations rather than
against the station records themselves. Each agrees with the
independent BMKG network far less: $0.195$ for IMERG-L and $0.154$
for CPC-UNI. That a gauge analysis assembled from station reports
should agree with an independent station network no better than a
raw satellite retrieval does is not a credible statement about data
quality; it points instead at the date label, which the two gridded
archives share and the BMKG network, reporting on the local
morning-observation day, does not.

\subsection{Agreement After Harmonising the Date Label}

Shifting each gridded series by whole days against the BMKG date
resolves the discrepancy. Every product peaks at a shift of
$-1$~day, the offset implied by the morning-observation convention:
the BMKG total dated~$d$ accumulates over the 24~hours ending at
the morning observation on~$d$ and therefore covers mostly rain
that fell on the previous calendar day. At that shift the
correlation rises to $0.751$ for CPC-UNI, a factor of $4.9$, and to
$0.568$ for IMERG-L, a factor of $2.9$, for a change that alters no
precipitation value.

The recovery is specific to the GPM era, and the reason is
instructive. Splitting the record at the arrival of the full GPM
constellation, the pre-GPM years of 2001 to 2013 give $0.358$ for
IMERG-L at the archived label against $0.320$ at the $-1$~day
label, a difference of little consequence, whereas the GPM years
give $0.195$ against $0.568$. The earlier product does not resolve
the diurnal cycle, so its daily total is temporally smeared and
correlates moderately whichever 24-hour window is used; the later
product resolves timing sharply and therefore correlates strongly
at the matched window and poorly at the mismatched one. Improved
retrieval timing makes a product more sensitive to a calendar
error, not less. The consequence for reporting is that the full
record sits between the two eras at both labels, so a correlation
against BMKG quoted without its era is ambiguous by roughly
$0.10$. This day-resolution result reproduces the flat pre-GPM
behaviour that the half-hourly diagnostic of
Section~\ref{sec:results-window} finds independently.

This day-label scan and the half-hourly window diagnostic of
Section~\ref{sec:results-window} are independent measurements and
they agree closely. The daily IMERG-L archive at zero shift gives
$0.1951$, against the $0.1949$ obtained by re-accumulating the
half-hourly source at $h = 0$; at the matched label the two give
$0.568$ and $0.571$. The small residual is expected, since the day
scan resolves no finer than $-24$~h while the half-hourly scan
places the optimum at $-23$~h, and the two pool marginally different
station sets. CPC-UNI evaluated on its native $0.5^{\circ}$ grid and
on the regridded $0.1^{\circ}$ grid agree to within $0.003$
($0.751$ and $0.749$ at the $-1$~day label), so the regridding step
carried out inside the pipeline does not affect this comparison.

The third pair completes the matrix from the other direction.
Table~\ref{tab:pairwise-daylabel} scans the two satellite-versus-gauge
pairs across date labels. The IMERG-L against CPC-UNI pair is scanned
instead across accumulation windows, because both archives are gridded
and neither carries a station reporting convention to shift; that scan
is reported in Chapter~\ref{chap:discussion}
(Table~\ref{tab:convention}), and it completes the diagonal: each
reference peaks at its own window and falls away at the other.

\subsection{What the Harmonised Values Do and Do Not Show}

At the $-1$~day label CPC-UNI reaches $0.751$ against BMKG while
IMERG-L reaches $0.568$. This gap is not a skill comparison.
CPC-UNI is a gauge analysis assembled from station reports
distributed over the Global Telecommunication System, and, as
Section~\ref{sec:methods-reference-quality} sets out, the
Indonesian contributions to that stream originate from BMKG. At the matched label the
reference is therefore partly reproducing the observations against
which it is being scored, and the per-station median at that label,
$0.865$, sits above the pooled value in a way that is consistent
with that shared origin.

Two conclusions follow. First, the low correlation between CPC-UNI
and BMKG under the archived labels is a calendar artefact rather
than a defect in the reference: harmonise the label and the two
gauge datasets agree closely. Second, and for the same reason,
CPC-UNI cannot serve as an independent check on a product
calibrated against it. The out-of-sample validation against BMKG
reported through the rest of this chapter is a requirement of the
design, not a supplementary test.

\subsection{The Same Comparison After Correction}

Repeating the scan on the three corrected stages separates the
calendar effect from the correction. Under the same construction
the correlation against BMKG moves from $0.195$ for the raw archive
to $0.198$, $0.173$, and $0.174$ for LS, LSEQM, and LSEQM+DL under
the archived label, and from $0.568$ to $0.569$, $0.552$, and
$0.555$ at the $-1$~day label (Table~\ref{tab:pairwise-daylabel}).
Two points follow.

No stage materially improves day-to-day timing at either label. LS
sits $0.003$ above the raw archive at the archived label and $0.001$
above it at the $-1$~day label, differences that are at the edge of
the sampling uncertainty for this sample size, and the two later
stages fall slightly below the raw archive at both. This closes an
objection that the archived-label numbers alone cannot answer: that
the flat correlation reported against the independent stations
might be an artefact of misaligned dates concealing a real gain. It
is not. With the dates aligned the corrected product still does not
separate from the raw archive. The corresponding in-sample result
against CPC-UNI is reported separately in
Section~\ref{sec:results-correlation} and
Table~\ref{tab:headline-cpc}, where the correction is likewise flat
but the values sit near $0.348$, being measured per pixel against the
calibration target rather than pooled over station-days at the gauges.

The calendar effect is independent of the correction. Every stage
retains the lift at $-1$~day, so the relabelling can be applied
before or after the pipeline for the same benefit. A separate
comparison supports the same reading: measured per pixel and
averaged over the $36$ dekadal windows of the full record, the
corrected fields correlate with the raw IMERG-F Final run at
$0.955$, $0.935$, and $0.936$ for the three stages, so the
correction rescales the field without restructuring its day-to-day
pattern.

\begin{table}[H]
\centering
\caption[Daily correlation against BMKG under two date labels]{Daily
Pearson $r$ against the $177$ BMKG stations, as a single pooled
correlation over the continuous record, under the archived date
label and after a $-1$~day shift; $n$ is $2.9\times10^{5}$ paired
station-days in the GPM era and $8.7\times10^{5}$ over the full
record.\label{tab:pairwise-daylabel}}
\begingroup\small
\begin{tabular*}{\textwidth}{@{\extracolsep{\fill}}lcccc@{}}
\toprule
& \multicolumn{2}{c}{GPM era, 2015 to 2021}
& \multicolumn{2}{c}{Full record, 2001 to 2021} \\
\cmidrule(lr){2-3}\cmidrule(lr){4-5}
Product & Archived & $-1$~day & Archived & $-1$~day \\
\midrule
IMERG-L (raw)   & 0.195 & 0.568 & 0.295 & 0.405 \\
LS              & 0.198 & 0.569 & 0.299 & 0.415 \\
LSEQM           & 0.173 & 0.552 & 0.276 & 0.403 \\
LSEQM+DL        & 0.174 & 0.555 & 0.279 & 0.406 \\
\addlinespace
CPC-UNI         & 0.154 & 0.751 & 0.332 & 0.532 \\
\bottomrule
\end{tabular*}
\endgroup
\end{table}


% ============================================================
\section{Stage-wise Headline Skill}\label{sec:results-headline}
% ============================================================

Each correction stage was designed to address a distinct bias
dimension, and the verification organises around the same
taxonomy. Ten metrics in Table~\ref{tab:headline} group into the
four pillars introduced in Section~\ref{sec:methods-metrics}:
\textit{value adjustment} (relative bias (RB), standard deviation
ratio (SDR)), \textit{distribution alignment} (Kolmogorov-Smirnov (KS) $p$-value, wet-day
frequency, wet-day intensity), \textit{extreme preservation}
($Q_{95}$ and $Q_{99}$ ratios), and \textit{event detection}
(Probability of Detection (POD), Critical Success Index (CSI), False Alarm Ratio (FAR)). Reading down a column gives the per-stage
performance; reading across pillars gives the design coverage of
the framework.

The table is reported stage-by-stage rather than
only as raw-versus-corrected so that each column can be audited
against the bias dimension its stage was designed to address. If
LS over-corrects the mean, the LS column will reveal it; if EQM
under-aligns the bulk distribution, the LSEQM column will too;
and if the CNN distorts a metric outside its operating range,
the LSEQM+DL column will show that as well. All numbers are
pooled across the 36 dekadal windows of 2001 to 2021 at the 172
BMKG stations not used in any correction step, so the comparison
tests whether the corrections generalise beyond the Climate
Prediction Center (CPC) Unified Gauge-Based Analysis of Global
Daily Precipitation (CPC-UNI) training reference.

\begin{table}[H]
\centering
\caption[Stage-wise skill against 172 independent BMKG stations]{Stage-wise skill against 172 independent BMKG stations,
pooled across all 36 dekads (2001 to 2021)\label{tab:headline}}
\begingroup\small
\begin{tabular}{@{}llcccc@{}}
\toprule
Pillar & Metric & LS & LSEQM & LSEQM+DL & Target \\
\midrule
Value Adjustment
  & Relative Bias        & $-$0.114 & 0.009 & \textbf{$-$0.006} & 0 \\
  & SDR                  & 0.71 & 1.03 & \textbf{1.00} & 1.0 \\
Distribution Alignment
  & KS $p$-value (\%)    & 0.01 & \textbf{19.07} & \textbf{19.07} & 100 \\
  & Wet-Day Freq.\ Ratio & 1.21 & \textbf{0.96} & 0.95 & 1.0 \\
  & Wet-Day Int.\ Ratio  & 0.71 & 1.06 & \textbf{1.05} & 1.0 \\
Extreme Preservation
  & $Q_{95}$ Ratio       & 0.74 & 1.07 & \textbf{1.05} & 1.0 \\
  & $Q_{99}$ Ratio       & 0.71 & 1.05 & \textbf{1.01} & 1.0 \\
Event Detection
  & POD                  & \textbf{0.78} & 0.65 & 0.65 & 1.0 \\
  & CSI                  & \textbf{0.53} & 0.49 & 0.49 & 1.0 \\
  & FAR                  & 0.36 & \textbf{0.32} & \textbf{0.32} & 0 \\
\bottomrule
\end{tabular}
\endgroup
\end{table}

Table~\ref{tab:headline} reports the out-of-sample skill against
the 172 BMKG validation stations, which is the strongest claim
the framework supports because the BMKG network is independent
of the correction reference. Bold values mark the best
performance per row, and the target column gives the value for
perfect agreement with observations. The parallel paper
\citep{ref:istanto2026mdpi}\MDPIstatus reports a companion table
of in-sample skill against the CPC-UNI calibration target;
Table~\ref{tab:headline-cpc} reproduces those values here so the
two-tier validation framework of
Section~\ref{sec:methods-validation} can be read in full. In its
temporal-skill rows, RMSE is the root mean square error and NSE
is the Nash-Sutcliffe efficiency.

\begin{table}[H]
\centering
\caption[In-sample skill against the CPC-UNI calibration target]{In-sample skill against the CPC-UNI calibration target,
pooled across all 36 dekads (2001 to 2025)\label{tab:headline-cpc}}
\begingroup\small
\begin{tabular}{@{}llcccc@{}}
\toprule
Pillar & Metric & LS & LSEQM & LSEQM+DL & Target \\
\midrule
Value Adjustment
  & Relative Bias        & \textbf{0.000} & 0.070 & 0.065 & 0 \\
  & SDR                  & \textbf{0.971} & 1.159 & 1.148 & 1.0 \\
Distribution Alignment
  & KS $p$-value (\%)    & \textbf{6.56} & 0.00 & 0.00 & 100 \\
  & Wet-Day Freq.\ Ratio & \textbf{1.040} & 0.947 & 0.946 & 1.0 \\
  & Wet-Day Int.\ Ratio  & \textbf{0.961} & 1.157 & 1.152 & 1.0 \\
Extreme Preservation
  & $Q_{95}$ Ratio       & \textbf{0.946} & 1.171 & 1.163 & 1.0 \\
  & $Q_{99}$ Ratio       & \textbf{0.967} & 1.212 & 1.196 & 1.0 \\
Event Detection
  & POD                  & \textbf{0.767} & 0.707 & 0.707 & 1.0 \\
  & CSI                  & \textbf{0.589} & 0.572 & 0.572 & 1.0 \\
  & FAR                  & 0.271 & 0.248 & \textbf{0.248} & 0 \\
\midrule
Temporal Skill
  & Pearson $r$          & 0.343 & 0.345 & \textbf{0.348} & 1.0 \\
  & RMSE (mm/day)        & \textbf{13.10} & 14.18 & 14.07 & 0 \\
  & NSE                  & \textbf{$-$0.273} & $-$0.548 & $-$0.524 & 1.0 \\
\bottomrule
\end{tabular}
\endgroup
\end{table}

The two tables tell a complementary story. Against CPC-UNI
(Table~\ref{tab:headline-cpc}), LS has zero relative bias by
construction and the LSEQM/LSEQM+DL stages overshoot the
reference distribution slightly at the upper percentiles. Against
BMKG (Table~\ref{tab:headline}), LS retains a residual bias of
$-11.4\%$, dominated by the limit of a single multiplicative
correction that does not adjust distribution shape, with a smaller
contribution from the gauge-to-gauge offset between CPC-UNI and
BMKG; the stage-wise attribution below confirms the correction-side
contribution dominates. The LSEQM/LSEQM+DL stages move the corrected
product close to unity at every pillar. The apparent overshoot
of LSEQM+DL Q$_{99}$ ratio against CPC-UNI ($1.20$) becomes a
near-perfect match against BMKG ($1.01$): the framework was
calibrated to CPC-UNI but reproduces BMKG more accurately because
CPC-UNI itself under-catches heavy rain. The Temporal Skill
section of Table~\ref{tab:headline-cpc} is the basis for the
Pearson correlation analysis in
Section~\ref{sec:results-correlation}.

Reading across the four pillars against the independent BMKG
stations (Table~\ref{tab:headline}), the framework moves three
pillars cleanly toward the gauge target and produces a designed
trade-off in the fourth. Every skill value quoted in the four
paragraphs that follow is measured against BMKG. The value-adjustment pillar shows
LSEQM and LSEQM+DL essentially eliminating the LS residual:
relative bias moves from $-11.4\%$ under LS to $+0.9\%$ under
LSEQM and $-0.6\%$ under LSEQM+DL, leaving the corrected product
within one percentage point of the gauge target. The standard
deviation ratio (SDR) tracks the same path, from $0.71$ (a $29\%$
spread compression under LS) through $1.03$ (slight over-spread
under LSEQM) to exactly $1.00$ under LSEQM+DL, where the CNN
refinement pulls the EQM overshoot back to the gauge target.

The distribution-alignment pillar shifts from rejection at the
$99.99\%$ level under LS ($p = 0.01\%$) to acceptance at
conventional confidence levels under LSEQM and LSEQM+DL ($p =
19.1\%$), meaning the corrected daily distributions are no
longer statistically distinguishable from the gauge reference.
The wet-day frequency ratio echoes the move, from $1.21$ (a
$21\%$ over-count of rainy days under LS, consistent with the
Integrated Multi-satellitE Retrievals for Global Precipitation Measurement (IMERG) tendency to over-detect drizzle and light rain) to $0.96$
and $0.95$ under LSEQM and LSEQM+DL respectively, indicating
that the dry-day-matching mechanism has done its work without
introducing a counter-bias. The wet-day intensity ratio moves
in parallel from $0.71$ under LS to $1.06$ under LSEQM and
$1.05$ under LSEQM+DL, confirming that the mean rainfall on wet
days now matches the gauge magnitude after the bulk-distribution
correction.

The extreme-preservation pillar recovers the upper tail that LS
truncates by construction. $Q_{99}$ moves from $0.71$ (a $29\%$
under-prediction of heavy-rain magnitude) through $1.05$ under
LSEQM (a slight overshoot from the EQM tail extrapolation) to
$1.01$ after CNN refinement, and $Q_{95}$ follows the same arc
($0.74 \rightarrow 1.07 \rightarrow 1.05$). The CNN stage acts
as a tail trim, pulling the EQM overshoot back toward unity
rather than amplifying it. This refinement is confined to a small
extreme-tail footprint, modifying under a fifth of land pixel-days
and predominantly downward, quantified in
Figure~\ref{fig:cnn-footprint} below.

The event-detection pillar reveals the expected trade-off.
Matching the dry-day frequency suppresses spurious light-rain
events, which lowers aggregate POD from $0.78$ to $0.65$ and
CSI from $0.53$ to $0.49$ but also reduces FAR from $0.36$ to
$0.32$. The aggregate numbers in
Table~\ref{tab:headline} mask a threshold-dependent reversal
that Section~\ref{sec:results-threshold} details: skill loss
is concentrated at light-rain thresholds where false alarms
drive the count, while skill at heavy-rain thresholds improves
substantially. Figure~\ref{fig:headline-bars} presents one
representative metric per pillar in bar form, with LS, LSEQM,
and LSEQM+DL side-by-side and the perfect-agreement target
marked.

\begin{figure}[H]
\centering
\includegraphics[width=\figFourOneWidth\textwidth]{fig_thesis_04_headline_bars.png}
\caption[Stage-wise skill at 172 BMKG stations for four representative metrics]{Stage-wise skill at 172 independent BMKG stations for
four representative metrics drawn from
Table~\ref{tab:headline}\label{fig:headline-bars}}
\end{figure}

A single-day spatial example of all six products on $1$~January~$2025$,
including the raw IMERG-L, the three correction stages, the IMERG
Final Run (IMERG-F) independent benchmark, and the CPC-UNI gauge
reference, is in
Appendix~\ref{app:atlas}~Figure~\ref{fig:atlas-six}.


% ============================================================
\section{Stage-wise Attribution}\label{sec:results-stages}
% ============================================================

The three correction stages can be read as an additive sequence,
each addressing a residual that the previous stage cannot reach.
This section walks through the stages in order and reports, for
each, the evaluation metrics that shift and the magnitude of the
shift, before the threshold-stratified detection analysis in
Section~\ref{sec:results-threshold} unpacks the apparent
event-detection regression seen in the aggregate.

\subsection*{What Linear Scaling Achieves}

Linear Scaling is the simplest correction in the pipeline, and
its performance against BMKG reflects what a single
mean-rescaling step can and cannot do. The relative bias against
BMKG of the LS-corrected field is $-11.4\%$. LS was designed to
match the dekadal mean against CPC-UNI, not against BMKG, so its
BMKG-side residual reflects two contributions: the limited
correction (a single multiplicative factor that does not adjust
distribution shape) and the gauge-to-gauge offset between the
CPC-UNI gridded analysis and BMKG point measurements
\citep{ref:chen2008}. The decomposition between these two
contributions can be inferred from the next subsection: the
LSEQM stage uses the same CPC-UNI training and brings the BMKG
residual to $+0.9\%$, indicating that the correction-side
contribution dominates and that the residual is not a pure
reference offset.

The standard deviation ratio against BMKG sits at $0.71$,
indicating that the corrected satellite field still varies less
than the gauge observations. The distribution shape is
essentially unchanged: the Kolmogorov-Smirnov test against the
BMKG distribution returns $p = 0.01\%$, far below conventional
acceptance thresholds. The wet-day frequency ratio of $1.21$
confirms that LS retains the IMERG Late Run (IMERG-L) tendency to over-detect
light rain. LS also leaves the upper tail truncated: $Q_{99}$
stands at $0.71$, indicating that the satellite under-predicts
the heaviest events by a factor of about $0.71$. The qualitative
picture is consistent with what LS was designed for: it adjusts
the mean of the dekadal field but does not modify day-by-day
distributional structure. Stages 2 and 3 of the pipeline are
introduced precisely to address this gap.

\subsection*{What LSEQM Adds: the Largest Single Jump}

The LS-to-LSEQM transition produces the largest single
improvement in the framework across all three pillars except
event detection. The standard deviation ratio moves from $0.71$
to $1.03$, correcting a $29\%$ underestimation of observed
variability. The Kolmogorov-Smirnov $p$-value rises from
$0.01\%$ to $19.1\%$, so the corrected and BMKG distributions
can no longer be rejected as different populations at
conventional confidence levels \cite{ref:ebert2007, ref:wilks2020}. The wet-day frequency ratio drops from $1.21$ to
$0.96$, removing nearly all of the $21\%$ overestimation of
wet-day occurrence through the dry-day frequency matching
scheme. The $Q_{99}$ ratio improves from $0.71$ to $1.05$ and
the $Q_{95}$ ratio from $0.74$ to $1.07$, recovering the heavy
tail that LS truncates by construction. The relative bias
against BMKG flips sign and shrinks to $+0.9\%$, because gamma
EQM rescales the full distribution rather than only the mean.

Two costs accompany these gains. The probability of detection
(POD) drops from $0.78$ to $0.65$ and the critical success
index (CSI) from $0.53$ to $0.49$. Both follow from the dry-day
frequency matching design: when the corrected wet-day frequency
falls toward unity, a fraction of IMERG-L light-rain events
previously counted as hits are reclassified as dry. The
reclassification reduces over-detection and improves the false
alarm ratio (FAR) from $0.36$ to $0.32$, at the price of a
lower hit rate at low thresholds \citep{ref:maggioni2016}. The
trade-off is the intended outcome of the design choice;
Section~\ref{sec:results-threshold} shows where it is
favourable and where it is not.

A complementary reading of the same contingency table answers a
question operational users ask directly: on a day the gauge
records as dry, does the corrected product also report dry? Dry
days, those below the $1$~mm wet-day threshold, account for
$50\%$ of the validated station-days, so their treatment carries as
much weight as wet-day skill. Pooling the $1$~mm contingency table
over all station-dekads rather than averaging per-station rates,
the fraction of dry gauge days that the product
also reports dry rises from $56.2\%$ under Linear Scaling to
$69.3\%$ under LSEQM+DL, so the false-alarm rate on dry days
falls from $43.8\%$ to $30.7\%$. This is the dry-day counterpart
of the wet-day frequency ratio moving from $1.21$ to $0.95$: the
empirical quantile mapping reclassifies spurious light-rain pixels
to dry and removes the raw product's over-detection. The gain is
not free: it is the dry-day face of the same reclassification that
costs light-rain hits above.

\subsection*{What LSEQM+DL Adds: Tail Refinement}

The contribution of Stages~3 and 4 over the LSEQM baseline is
concentrated in the extreme tail. The $Q_{99}$ ratio sharpens
from $1.05$ to $1.01$, pulling the upper tail to within $1\%$
of the BMKG target; the $Q_{95}$ ratio tightens from $1.07$ to
$1.05$; and the SDR moves from $1.03$ to exactly $1.00$.
Distribution-alignment metrics (KS $p$-value, wet-day frequency
and intensity ratios) and event-detection metrics (POD, CSI,
FAR) are essentially unchanged because the CNN refinement
operates only above the $80$th-percentile threshold and does not
touch the bulk of the distribution.

The spatial footprint of the CNN confirms that it is a targeted
refinement rather than a whole-field rewrite.
Figure~\ref{fig:cnn-footprint} shows the share of land pixel-days
the CNN modifies in each year: on average $18.6\%$ over the
$2001$ to $2025$ record, confined to the pixels above the
$80$th-percentile tail. About two-thirds of the modified pixels
($67\%$) are adjusted downward, trimming the extreme overshoot the
Gamma and GPD extrapolation leaves behind, with a mean absolute
change of $0.36$~mm/day. These individually small, spatially
concentrated adjustments remove about four-fifths of the residual
$99$th-percentile overshoot while leaving the day-by-day pairing,
and hence Pearson~$r$, pinned close to its raw value
(Section~\ref{sec:discussion-bound}). The CNN thus acts only where
the statistical stages leave an incoherent residual in the
extremes, which is why it sits last in the pipeline.

\begin{figure}[H]
\centering
\includegraphics[width=\figFourCnnWidth\textwidth]{fig_thesis_I1_cnn_footprint.png}
\caption[CNN refinement footprint by year]{Share of land pixel-days the CNN modifies each year, with the whole-record mean; the CNN acts on a small extreme-tail footprint.\label{fig:cnn-footprint}}
\end{figure}

The stage-resolved attribution is therefore
complete: each stage moves the metrics matched to its design
dimension and leaves the others alone.
Appendix~\ref{app:atlas}~Figure~\ref{fig:atlas-diff} maps the
per-pixel residual on the same representative day, locating
where in the archipelago the correction is doing the most work.


% ============================================================
\section{Threshold-Dependent Detection}\label{sec:results-threshold}
% ============================================================

Daily precipitation over the region is strongly right-skewed, and
the correction has to act on that shape rather than on a
symmetric error. This section is reported against BMKG only: the
threshold-stratified contingency tables are accumulated at the
station points, and no gridded CPC-UNI equivalent exists.

Pooled over the daily BMKG station records at the satellite pixels
containing them, $2001$ to $2021$, wet-day
amounts have a skewness of $3.2$, the wettest $5\%$ of wet days
carry $26\%$ of the total and the wettest $1\%$ carry $8\%$, and
the $99$th wet-day percentile is $106$~mm/day against a
wet-day mean of $17$~mm/day. A heavy upper tail of this kind is
what the Generalized Pareto step of Stage~3 is included to
preserve; a Gamma fit alone underestimates the far tail of such a
distribution. The raw IMERG-L input departs from this reference
shape in two directions at once: it reports too many wet days (a
wet-day fraction of $0.61$ against $0.46$ for the gauges) while
truncating the heavy tail (a $99$th percentile of $73$~mm/day
against $106$). These are the two defects the threshold-stratified
verification exposes, and they move detection skill in opposite
directions as the threshold rises.

The aggregate event-detection metrics in
Table~\ref{tab:headline} mask a threshold-dependent reversal in
detection skill. Skill is broken out at a series of daily
thresholds in Figure~\ref{fig:threshold-skill}. Thresholds up to
$50$~mm/day follow the World Meteorological Organization (WMO)
multi-threshold verification practice \citep{ref:wmo1485,
ref:wmo1132}, while $100$ and $150$~mm/day are added as heavy and
extreme classes relevant to the Indonesian wet-season regime,
consistent with the WMO principle that extreme-precipitation
thresholds are set regionally \citep{ref:wmo1310}. The thresholds
are labelled here by their conventional daily-rainfall intensity
descriptors: $1$~mm (measurable), $5$~mm (light), $10$~mm
(moderate), $20$~mm (moderate-heavy), $50$~mm (heavy), and
$100$~mm (very heavy). The $150$~mm/day extreme class is omitted
because no station reaches the minimum of ten exceedances required
for stable categorical verification in any dekad of the record.

The two input defects account for the reversal. At the
$1$, $5$, and $10~\mathrm{mm/day}$ thresholds LS retains a higher
POD than LSEQM+DL. LS is a single multiplicative factor that leaves
the wet-day count untouched, so it keeps the raw over-detection of
light rain, and an over-detecting product scores more hits at a low
threshold. The distribution step of LSEQM+DL removes that
over-detection, reclassifying the spurious drizzle to dry, which
costs some of those low-threshold hits. The ordering changes only at
$20~\mathrm{mm/day}$, not at $5~\mathrm{mm/day}$.

Above the crossover the same two edits act the other way. LS
inherits the truncated tail, so it under-detects the heavy events
that sit above $20~\mathrm{mm/day}$. The empirical quantile
mapping and Generalized Pareto steps, by contrast, restore the upper tail from a
$99$th percentile of $73$ toward the gauge value of $106$~mm/day,
recovering heavy events the raw product places below the threshold.
The crossover between $10$ and $20~\mathrm{mm/day}$ is where the
loss of over-detected light rain and the recovery of the truncated
tail balance, and from $20~\mathrm{mm/day}$ upward LSEQM+DL
achieves the higher CSI at every resolved threshold.

At $50~\mathrm{mm/day}$ the median CSI is $0.091$ for LSEQM+DL
against $0.062$ for LS, a $45\%$ improvement, and the equitable
threat score (ETS), which removes the chance-hit contribution,
shows the same ordering. These two values are medians taken over
the $983$ station-dekad records that meet the ten-exceedance
minimum; Figure~\ref{fig:threshold-skill} instead averages the
per-dekad station medians, which raises the LSEQM+DL value to
$0.097$ and the improvement to $56\%$ without changing the
ordering. The $100~\mathrm{mm/day}$ class supports
no ordering: only $6$ of the $5{,}899$ station-dekad records reach
the ten-exceedance minimum, and the median CSI is zero for both
methods. The advantage is therefore established at $20$ and
$50~\mathrm{mm/day}$, where the record carries enough events.

The crossover threshold near $20~\mathrm{mm/day}$ is not an
arbitrary cut. It is the fixed threshold of R20mm, the
very-heavy-precipitation-day index of the World Meteorological
Organization Expert Team on Climate Change Detection and Indices,
defined as the annual count of days with precipitation at or above
$20~\mathrm{mm}$; the corresponding heavy-precipitation index
R10mm sits at $10~\mathrm{mm}$ and brackets the crossover from
below \citep{ref:peterson2001, ref:zhang2011}. LSEQM+DL therefore
overtakes LS in the regime that standardised indices class as very
heavy precipitation, which is what threshold-based warning over the
Maritime Continent depends on \citep{ref:ramadhan2022}.

\begin{figure}[H]
\centering
\includegraphics[width=\figFourTwoWidth\textwidth]{fig_thesis_04_multi_threshold.png}
\caption[Threshold-stratified verification at the $172$ BMKG station locations]{Detection skill as a function of daily rainfall threshold at the $172$ BMKG
station locations\label{fig:threshold-skill}}
\end{figure}


% ============================================================
\section{Pearson Correlation Across Stages}\label{sec:results-correlation}
% ============================================================

Section~\ref{sec:results-intercomparison} established, against the
independent BMKG stations and at both date labels, that the
correction does not lift the daily correlation. This section
reports the same result against the in-sample CPC-UNI target, so
the finding holds against both references.

One metric does not move across the three corrected stages. The Pearson
correlation between paired daily satellite and CPC-UNI values,
computed at each pixel and pooled across the dekad, is $0.343$
for LS, $0.345$ for LSEQM, and $0.348$ for LSEQM+DL. The total
shift across stages is approximately $+0.005$, well within the
single-pixel sampling uncertainty of $\pm 0.06$ implied by the
Fisher $z$-transform standard error at $n = 250$
\citep{ref:wilks2020}.

Figure~\ref{fig:paradox} states the central observation of this
chapter in four panels. Pearson $r$, root mean square error
(RMSE), and Nash-Sutcliffe efficiency (NSE)
\citep{ref:nashsutcliffe1970} are temporal-skill diagnostics
measured here against the in-sample CPC-UNI target, and they do
not improve across the three corrected stages: $r$ stays near
$0.35$ across all three, RMSE rises from $13.1$ to $14.1$~mm/day
because the correction widens spread without re-aligning timing,
and NSE moves from $-0.27$ to $-0.52$. The standard-deviation
ratio, an amplitude diagnostic measured against the independent
BMKG gauges, moves from $0.71$ to $1.00$ across the same stages.
The amplitude is corrected to the validation target, yet the
timing does not improve even against the reference the correction
was fitted to.

\begin{figure}[H]
\centering
\includegraphics[width=\figFourThreeWidth\textwidth]{fig_thesis_04_paradox.png}
\caption[What does and does not improve under marginal bias correction]{What does and does not improve under marginal bias
correction, with the temporal-skill panels scored against CPC-UNI and the amplitude panel against BMKG.\label{fig:paradox}}
\end{figure}

The contrast is not a verification quirk; it is what the
framework should do given how each metric is built. The SDR
depends only on the marginal variance of the two series and is
therefore moved directly by the EQM stage acting on the daily
distribution. Pearson $r$, RMSE, and NSE depend on the
day-by-day pairing of satellite and gauge values and inherit
whatever timing offset the satellite product carries from its
half-hourly inputs and the daily aggregation window. The
marginal correction has no mechanism to shift a value from one
day to its neighbour, so by construction it cannot lift the
three temporal-skill metrics. The visual asymmetry between the
amplitude-pillar improvement and the temporal-skill diagnostics
is the central observation of this thesis;
Section~\ref{sec:discussion-bound} develops the theoretical
basis for the bound.

The empirical value $r \approx 0.348$ is of the
order reported for daily IMERG over the
Indonesian Maritime Continent \citep{ref:ramadhan2022ts}, and well
below the median near $0.75$ that the same product reaches against
the far denser Chinese gauge analysis
\citep[Fig.~3]{ref:tang2020}. That gap is not left unexplained:
Section~\ref{sec:results-window} shows most of it is a
calendar-window effect rather than a retrieval one.

The single figure $r \approx 0.35$ is the spatial mean of a broad
per-pixel distribution. Table~\ref{tab:cpc-corr-dist} and
Figure~\ref{fig:cpc-corr-dist} characterise the daily correlation
against CPC-UNI across the full $0.1^{\circ}$ domain: over the
$19{,}393$ land pixels it spans $0.08$ to $0.68$, with a median of
$0.35$ and a pooled value of $0.36$.

The correlation is markedly
higher where the CPC-UNI gauge network is dense: at the $169$ cells
that host at least one of the $180$ BMKG stations the median and the
pooled value are both $0.51$, against $0.35$ and $0.36$
domain-wide. The gridded
reference is better constrained where the gauges cluster. At every aggregation the
value is flat across the three corrected stages (whole-domain
pooled $0.350$, $0.360$, and $0.363$), so the near-flat behaviour
of $r$ is a property of the whole field, not of a particular sample
of points. All of these are measured against CPC-UNI, the in-sample
calibration target; the independent, out-of-sample comparison
against the BMKG point gauges is lower and is decomposed in
Section~\ref{sec:results-window}, where the day-by-day pairing
against the station observations proves to be dominated by the
calendar-window mismatch rather than by retrieval skill.

\begin{table}[H]
\centering
\caption[Daily correlation against CPC-UNI across the domain]{Daily Pearson $r$ of the corrected product against CPC-UNI across the domain over the full 2001 to 2025 record, by correction stage\label{tab:cpc-corr-dist}}
\begingroup\small
\begin{tabular*}{\textwidth}{@{\extracolsep{\fill}}lccccc@{\hskip 1.4em}c@{}}
\toprule
& \multicolumn{5}{c}{Whole coverage ($19{,}393$ pixels)} & Gauge cells \\
\cmidrule(lr){2-6}
Stage & Min & Mean & Median & Max & Pooled & Pooled \\
\midrule
LS       & $0.092$ & $0.351$ & $0.341$ & $0.670$ & $0.350$ & $0.503$ \\
LSEQM    & $0.082$ & $0.352$ & $0.344$ & $0.673$ & $0.360$ & $0.508$ \\
LSEQM+DL & $0.082$ & $0.355$ & $0.346$ & $0.677$ & $0.363$ & $0.513$ \\
\bottomrule
\end{tabular*}
\endgroup
\end{table}

Figure~\ref{fig:cpc-corr-dist} renders the same whole-coverage
distribution as a density and overlays the gauge-dense cells. The
domain forms a single broad mode centred near $0.35$, while the
gauge-dense cells sit as a distinct, right-shifted population
centred near $0.51$, so the density effect is visible directly
rather than only through the pooled summary.

\begin{figure}[H]
\centering
\includegraphics[width=0.98\textwidth]{fig_thesis_04_cpc_corr_dist.png}
\caption[Spatial distribution of the CPC-UNI daily correlation]{Per-pixel daily Pearson $r$ of the LSEQM+DL product against CPC-UNI: the whole-coverage field against the gauge-dense cells\label{fig:cpc-corr-dist}}
\end{figure}


% ============================================================
\section{Temporal Alignment and the Optimal Aggregation
Window}\label{sec:results-window}
% ============================================================

The near-flat correlation of
Section~\ref{sec:results-correlation} raises a question the stage
comparison cannot answer: how much of the absolute ceiling is set
by the retrieval timing, and how much by the calendar window used
to label each daily total. The daily IMERG-L archive accumulates
over the Coordinated Universal Time (UTC) day, whereas the BMKG
stations report under the morning-observation convention, a
24-hour total ending near $07{:}00$ local time. To separate the two
components, the half-hourly IMERG-L source is re-aggregated into
24-hour windows starting at every integer-hour offset
$h \in [-48, +48]$ from $00{:}00$~UTC and correlated against the
BMKG daily totals (Section~\ref{sec:methods-window}). The
diagnostic operates on the raw retrieval rather than the corrected
product. The diagnostic pools the $178$ of the $180$
BMKG stations that carry valid half-hourly coverage, a different
subset from the $172$ stations used for the dekadal
correction-metric validation.

Figure~\ref{fig:window-era} reports the result over the GPM era,
$2015$ to $2021$. The pooled correlation rises from $r = 0.20$ at
$h = 0$, the UTC-day alignment the daily archive uses, to
$r = 0.57$ at $h = -23$, the window that matches the local-day
gauge convention. This UTC-day value of $0.20$ is lower than the
dekad-sliced CPC-UNI value of $r \approx 0.35$
(Section~\ref{sec:results-correlation}) for two reasons: the full
record blends the GPM era with the earlier, timing-imprecise era,
and the window diagnostic pairs against the independent BMKG gauges
rather than the CPC-UNI reference used in
Section~\ref{sec:results-correlation}. Each contributes about
$0.10$, so neither alone accounts for the difference.

The recovery is
conditional on the satellite era. Before the full GPM
constellation, the years through $2013$, whose IMERG inputs derive
from the sparser TRMM-era sensor set, re-windowing does not lift
the correlation: it stays near $r = 0.36$ at every offset and the
peak offset wanders without a stable value. The window-matched
recovery appears only once the half-hourly field resolves the
diurnal cycle, a transition the per-year trace places at the
$2014$ to $2015$ boundary.

\begin{figure}[H]
\centering
\includegraphics[width=\figFourNineWidth\textwidth]{fig_thesis_04_window_era.png}
\caption[Window-offset correlation recovery by satellite era]{Per-year
Pearson~$r$ at the UTC-day window ($h = 0$) and at each year's
best-matched window, with the peak offset~$h^{\star}$ below.\label{fig:window-era}}
\end{figure}

The optimal offset is stable across the seasonal cycle. Across all
twelve three-month running seasons the pooled peak sits at
$h = -23$~h in every timezone band
(Appendix~\ref{app:window}~Figure~\ref{fig:window-seasonal}), and
the optimum does not
migrate between the wet and dry monsoon phases.

The offset is equally stable across space. The per-station peak
offset has a median of $-23$~h, an interquartile range of
$[-23, -21]$~h, and no systematic longitudinal organisation across
the WIB, WITA, and WIT timezone bands
(Appendix~\ref{app:window}~Figure~\ref{fig:window-map}). The recovered per-station
correlation has a median of $0.56$, an interquartile range of
$0.51$ to $0.62$, and a maximum near $0.78$; taking the same
per-station median within each timezone band gives $0.55$ for WIB
and WITA and $0.61$ for WIT. The offset is therefore a property of the gauge
reporting convention shared across the network, not a local effect.

The window choice governs daily skill but not monthly skill.
Aggregating the same daily totals to monthly sums and validating
against the BMKG monthly totals returns a correlation near $0.80$
whether the daily windows are UTC-aligned or gauge-matched
(Table~\ref{tab:window-monthly}): a one-day window shift only
redistributes one or two days across each month boundary, so it is
invisible at the monthly scale. The same product that scores
$r = 0.20$ against the daily gauges at the UTC-day alignment scores
$r = 0.80$ against the monthly gauges, which places the daily
ceiling in the timing of the day-by-day pairing rather than in the
amount of rainfall the product captures.

\begin{table}[H]
\centering
\caption[Daily versus monthly correlation under two aggregation windows]{Daily
and monthly Pearson~$r$ under the UTC-day and gauge-matched windows
(GPM era; $n = 5{,}338$ station-months)}
\label{tab:window-monthly}
\begin{tabular*}{\textwidth}{@{\extracolsep{\fill}}lcc@{}}
\toprule
Aggregation window & Daily $r$ & Monthly $r$ \\
\midrule
UTC day ($h = 0$)        & $0.20$ & $0.80$ \\
Gauge-matched ($h = -23$) & $0.57$ & $0.81$ \\
\bottomrule
\end{tabular*}
\end{table}

These results localise the near-flat correlation of
Section~\ref{sec:results-correlation}: the larger part of the gap
between the UTC-day value and the gauge-matched value is a
calendar-window component, recoverable by re-aggregating the
archive over a window matched to the gauge convention, and the
remainder is a residual retrieval-skill component that no window
choice removes. The four-pillar value, distribution, extreme, and
detection scores are functions of the marginal distribution and of
categorical events, and are unchanged by the window observation.
Section~\ref{sec:discussion-subdaily} develops the interpretation:
why the recovery is confined to the GPM era, and what the
re-aggregation implies operationally.
Appendix~\ref{app:window} carries the supporting diagnostics for
this section, including the per-offset curves behind the summary
values reported here.


% ============================================================
\section{Spatial Skill Distribution}\label{sec:results-cqi}
% ============================================================

The headline numbers in Section~\ref{sec:results-headline} pool
across the full Indonesian domain and so average over the
regional heterogeneity of Java-Sumatra-Bali gauge density against
the sparse network in Maluku and Papua. The Composite Quality
Index (CQI) defined in Section~\ref{sec:methods-metrics}
collapses the per-pixel performance across the three sub-scores
(Basic Statistical Quality sub-score (BSQS), Distribution Quality sub-score (DQS), Temporal Quality sub-score (TQS)) and lets the spatial distribution of skill be
read as a single map. Figure~\ref{fig:cqi-spatial} shows the
CQI for the three correction stages, computed per pixel against
the in-sample CPC-UNI target.

\begin{figure}[H]
\centering
\includegraphics[width=\figFourFourWidth\textwidth]{fig_thesis_04_cqi_spatial.png}
\caption[CQI maps for the three correction stages]{CQI maps for (top) Linear Scaling, (middle) LSEQM, and (bottom) LSEQM+DL, computed
pixel-wise against CPC-UNI and averaged across all $36$ dekadal
windows.\label{fig:cqi-spatial}}
\end{figure}

The map view makes the regional heterogeneity visible: high
quality across Java and southern Sumatra where the gauge network
is dense, progressively lower quality eastward into Maluku and
Papua where the reference dataset itself is poorly constrained.
The differences between LSEQM and LSEQM+DL are concentrated in
the densely gauged regions, where the CNN refinement contributes
through the station-density-aware blending of
Section~\ref{sec:methods-blending}; in sparsely gauged regions
the framework reverts toward the LSEQM result by design, so the
LSEQM and LSEQM+DL maps are nearly indistinguishable there.

Appendix~\ref{app:atlas}~Figure~\ref{fig:cqi-improvement} isolates the LSEQM+DL
contribution by mapping the per-pixel CQI change from LSEQM to
LSEQM+DL (panel~a) and the LSEQM+DL CQI itself (panel~b) against
the four operational quality tiers (Poor $<0.4$, Fair $0.4$ to $0.6$,
Good $0.6$ to $0.8$, Excellent $\geq 0.8$). The change from the CNN
step is small and never negative. It stays within $\pm 0.01$ CQI
across most of the domain and rises to modest positive values over
the densely gauged western islands where the refinement is most
active (mean $+0.003$, exceeding $+0.01$ over $3.8\%$ of pixels and
falling below $-0.01$ over none).

The CQI map places the product in
the Fair band ($0.40$ to $0.60$) across $98.5\%$ of the land mass
(mean $0.503$), with a west-east gradient that tracks gauge
density. The product sits higher within the Fair band over Java, Sumatra, and the
well-observed interiors, and drops into the Poor tier over the
$1.5\%$ of pixels in the most sparsely gauged east. It does not
reach the Good or Excellent tier anywhere; against the sparse CPC
reference the attainable CQI is bounded within the Fair band,
consistent with the confidence mask reverting toward LSEQM where
stations are sparse.


The spatial pattern carries an important operational reading.
The headline LSEQM+DL numbers reported in
Table~\ref{tab:headline} apply with full strength in the
densely gauged west; in the eastern sparsely gauged regions the
relevant performance baseline is the LSEQM column rather than
the LSEQM+DL column, because the blending formula in
Section~\ref{sec:methods-blending} suppresses the CNN
contribution there. Users in Java, Sumatra, and Bali receive
the full benefit of the four-stage pipeline; users in Maluku
and Papua receive the LSEQM result with a thin and locally
variable CNN residual on top. The framework therefore does not
over-claim accuracy in regions where the reference dataset
itself cannot support refinement, which is the operational
expression of the design choice to make the CNN density-aware
rather than uniform. Pillar-by-pillar climatological metric maps
for each correction stage, on the same Indonesian domain, are in
Appendix~\ref{app:atlas}~Figure~\ref{fig:atlas-pillars}.


% ============================================================
\section{Per-Station Taylor Diagram Analysis}\label{sec:results-taylor}
% ============================================================

The aggregate pattern reported in
Sections~\ref{sec:results-headline} and
\ref{sec:results-correlation} can be cross-checked at the level
of individual BMKG stations using a Taylor diagram, which plots
each station as a point whose angular position encodes the
Pearson correlation against the gauge series and whose radial
position encodes the standard deviation ratio relative to the
gauge target. Figure~\ref{fig:taylor-station} shows the
station-level diagram broken out by season: the wet season
(November to April) and the dry season (May to October), which
are the operationally meaningful periods for Indonesian
agriculture and water management.

\begin{figure}[H]
\centering
\includegraphics[width=\figFourSixWidth\textwidth]{fig_thesis_04_taylor_by_station.png}
\caption[Station-level Taylor diagram of BMKG validation by season]{Station-level Taylor diagram of BMKG validation by
season\label{fig:taylor-station}}
\end{figure}

The Taylor diagram resolves the headline pattern across
individual stations: the three method symbols cluster at the
same angular position (correlation) while their radial positions
(standard deviation) differ substantially, with LSEQM and
LSEQM+DL sitting close to the reference radius and LS sitting
short. The pattern is conserved across both seasons, so the
flat-correlation observation is not driven by either the
convective wet season or the gauge-sparse dry season alone but
holds station-by-station and season-by-season across the
archipelago. This per-station visual evidence anticipates the
theoretical argument in Section~\ref{sec:discussion-bound}: the
correlation stays pinned close to the underlying retrieval timing,
which is a property of the satellite product at each pixel and
each day, and is not lifted by the marginal correction at any
of the $178$ stations carrying a valid correlation in this
diagnostic.
Appendix~\ref{app:atlas}~Figure~\ref{fig:atlas-stations}
places the validation stations directly on dekadal-mean
precipitation maps of the four products as a visual spatial
agreement check, and Appendix~\ref{app:stations} reports the
per-station statistics behind this diagram for a selected subset
of the network.


% ============================================================
\section{Temporal Stability Across Dekads}\label{sec:results-temporal}
% ============================================================

The dekadal pooling design of
Section~\ref{sec:methods-validation} aggregates each of the
$36$ dekadal windows into a single metric value, which masks
the within-year variability that drives operational use across
the Indonesian wet and dry seasons. This section recovers the
monthly view by re-aggregating the dekadal metrics to a
calendar-month axis and reports four representative metrics
(SDR, RMSE, KS $p$-value, CSI) across the three correction
stages, in Figure~\ref{fig:temporal-stability}. These are the
per-pixel metrics measured against the
in-sample CPC-UNI target, not the per-station BMKG values of
Section~\ref{sec:results-headline}.

\begin{figure}[H]
\centering
\includegraphics[width=\figFourSevenWidth\textwidth]{fig_thesis_04_temporal_stability.png}
\caption[Monthly temporal stability of four representative metrics across the three correction stages]{Monthly temporal stability of four representative
metrics across the three correction stages, measured per pixel against CPC-UNI\label{fig:temporal-stability}}
\end{figure}

The SDR stays near $1.0$ for LSEQM and LSEQM+DL across all
months, indicating that the spread correction is stable across
the seasonal cycle and does not require season-specific
calibration. RMSE rises during the peak wet season (December
to February) when absolute precipitation magnitudes are
largest, which is a property of any squared-error metric on a
distribution with a wet-season fat tail rather than a
correction artefact. The KS $p$-value remains above the
$5\%$ acceptance threshold for both LSEQM and LSEQM+DL across
all months except the deepest wet-season months, where the
high counts of heavy-rain days reduce the test power. CSI
shows a modest seasonal cycle, but the ordering seen in the
per-pixel panel, with LS above LSEQM and LSEQM+DL at the
$1$~mm wet-day threshold at which the gridded CSI is computed,
is conserved month-by-month. The threshold-dependent reversal
of Section~\ref{sec:results-threshold} is a BMKG station result
and is not visible here.
A daily-scale spatial example of the day-to-day variation across
three contrasting days in the same dekad is in
Appendix~\ref{app:atlas}~Figure~\ref{fig:atlas-threedays}.

The monthly view confirms that the per-pixel medians in
Table~\ref{tab:headline-cpc} are not driven by a small subset of
dekads. The pattern of stage-wise attribution is consistent across
the seasonal cycle: each stage moves the metrics matched to
its design dimension in every month, and the temporal-skill
ceiling on Pearson $r$ persists in every month. The dekadal
pooling design is therefore a legitimate way to aggregate
results across years for the comparison framework of
Section~\ref{sec:results-headline}, because the underlying
pattern does not change qualitatively from month to month.


% ============================================================
\section{Sensitivity Analyses}\label{sec:results-sensitivity}
% ============================================================

Three parameters in the pipeline were set by domain reasoning
rather than by a formal cross-validated search: the CNN blending
weight $\alpha = 0.70$, the GPD threshold at the $80$th wet-day
percentile, and the station-density saturation count
$c_{\text{sat}} = 2$. This section states the expected
direction of effect of each, drawn from the design reasoning in
Chapter~\ref{chap:methods}, then reports a parameter sweep over
the Bali reproducibility subdomain
(Section~\ref{sec:results-reproducibility}) that tests those
predictions empirically. The subdomain is used because each of the
fifteen settings requires a complete correction and
metric-evaluation run, including per-dekad CNN training: the sweep
is tractable on the $126$-pixel Bali grid, but the per-pixel stages
scale with the pixel count, and for the $19{,}395$ land
pixels of the full domain a fifteen-run national sweep is impractical
on commodity hardware. A full-domain sweep with a strict temporal
hold-out remains future work.

The blending weight $\alpha$ sets the maximum CNN contribution
in fully gauge-supported pixels through
Equation~\eqref{eq:alpha-eff}: a larger $\alpha$ moves the
blend toward pure LSEQM and reduces the CNN influence, while a
smaller $\alpha$ admits more of the learned residual. The
operating value $0.70$ caps the CNN at a $30\%$ contribution
where confidence is full, a deliberately conservative setting
that limits the spatial refinement to a correction rather than
a replacement.

The GPD threshold percentile trades sample size
against tail purity. Lowering it from the $80$th toward the
$70$th admits more exceedances and stabilises the maximum-likelihood
fit, but conflates the extreme regime with moderate rain that the
Gamma EQM already handles. Raising it toward the $90$th
sharpens the extreme regime at the cost of falling below the
exceedance count needed for a stable fit in dry pixels
(Section~\ref{sec:methods-gpd}).

The saturation count
$c_{\text{sat}}$ sets how many contributing gauges per CPC cell
are needed for full confidence. A larger value is more
conservative and suppresses the CNN across more of the domain,
while a smaller value admits the CNN more widely. The value $2$
was chosen after an initial setting of $3$ was found to suppress
the CNN almost everywhere outside Java.

The sweep varies each parameter through five values while
holding the other two at their operating defaults, runs the full
LSEQM+DL pipeline end-to-end on the Bali subdomain for all $36$
dekadal windows of the $2001$ to $2025$ record, and pools the
resulting per-pixel metrics.
Table~\ref{tab:sensitivity-bali} reports the fifteen settings,
with the operating default of each parameter shown in bold and
the per-pixel metrics given as $r$ (Pearson correlation), RB
(relative bias, \%), SDR (standard-deviation ratio), RMSE (root
mean square error, mm/day), and NSE (Nash-Sutcliffe efficiency).

\begin{table}[H]
\centering
\caption[Bali sensitivity sweep over the three configuration parameters]{Bali sensitivity sweep over the three configuration
parameters identified in Chapter~\ref{chap:methods}, reported per pixel against CPC-UNI at the native window over the $2001$ to $2025$ record\label{tab:sensitivity-bali}}
\begingroup\small
\begin{tabular*}{\textwidth}{@{\extracolsep{\fill}}lcccccc@{}}
\toprule
Parameter & Value & $r$ & RB (\%) & SDR & RMSE & NSE \\
\midrule
\multicolumn{7}{l}{\textit{Blending weight} $\alpha$} \\
$\alpha$ & $0.50$ & $0.344$ & $-14.5$ & $0.910$ & $5.88$ & $-0.40$ \\
         & $0.60$ & $0.343$ & $-14.9$ & $0.907$ & $5.89$ & $-0.40$ \\
         & $\mathbf{0.70}$ & $\mathbf{0.340}$ & $\mathbf{-15.0}$ & $\mathbf{0.905}$ & $\mathbf{5.89}$ & $\mathbf{-0.41}$ \\
         & $0.80$ & $0.336$ & $-15.2$ & $0.903$ & $5.90$ & $-0.41$ \\
         & $0.90$ & $0.332$ & $-15.5$ & $0.900$ & $5.91$ & $-0.41$ \\
\midrule
\multicolumn{7}{l}{\textit{GPD threshold percentile}} \\
Pctl     & $70$ & $0.345$ & $-13.0$ & $0.923$ & $5.96$ & $-0.42$ \\
         & $75$ & $0.342$ & $-13.7$ & $0.919$ & $5.96$ & $-0.41$ \\
         & $\mathbf{80}$ & $\mathbf{0.340}$ & $\mathbf{-15.0}$ & $\mathbf{0.905}$ & $\mathbf{5.89}$ & $\mathbf{-0.41}$ \\
         & $85$ & $0.337$ & $-16.2$ & $0.891$ & $5.82$ & $-0.40$ \\
         & $90$ & $0.335$ & $-18.0$ & $0.867$ & $5.72$ & $-0.39$ \\
\midrule
\multicolumn{7}{l}{\textit{Density saturation count} $c_{\text{sat}}$} \\
$c_{\text{sat}}$ & $1$ & $0.348$ & $-14.7$ & $0.911$ & $5.88$ & $-0.40$ \\
                 & $\mathbf{2}$ & $\mathbf{0.340}$ & $\mathbf{-15.0}$ & $\mathbf{0.905}$ & $\mathbf{5.89}$ & $\mathbf{-0.41}$ \\
                 & $3$ & $0.336$ & $-15.3$ & $0.903$ & $5.90$ & $-0.41$ \\
                 & $4$ & $0.334$ & $-15.3$ & $0.902$ & $5.90$ & $-0.41$ \\
                 & $5$ & $0.333$ & $-15.4$ & $0.901$ & $5.91$ & $-0.41$ \\
\bottomrule
\end{tabular*}
\endgroup
\end{table}

Three findings follow. First, every parameter produces a
monotone, locally smooth response across the explored range, and
none of the fifteen settings reverses the headline pattern of
Section~\ref{sec:results-headline} that LSEQM+DL moves value
adjustment, distribution alignment, and tail preservation toward
the gauge target.

Second, the Pearson correlation stays in the
narrow band $[0.332, 0.348]$ across all fifteen settings, a
spread of $0.016$. No setting shifts the daily correlation by
more than this amount, consistent with the interpretation of
Section~\ref{sec:results-correlation} that the daily correlation
is set by the timing skill of the raw retrieval rather than
by the choice of correction parameter \citep{ref:maraun2013}. This
flatness follows from the rank-preserving nature of the correction,
not from any property of the Bali subdomain, so it is a structural
result that holds wherever the pipeline is applied: the
correction-parameter insensitivity of the daily correlation is not
specific to Bali. The magnitudes of the distribution and bias
responses in Table~\ref{tab:sensitivity-bali}, by contrast, are
subdomain-specific and would differ in size over other regions, even
though their direction of effect follows the design reasoning above.

Third, the GPD threshold percentile dominates the dispersion:
raising the percentile from the $70$th to the $90$th narrows
the corrected distribution (SDR falls from $0.923$ to $0.867$)
and shifts relative bias more negative (from $-13.0\%$ to
$-18.0\%$), while RMSE falls modestly (from $5.96$ to
$5.72$~mm/day). The blending weight and the saturation count
each move bias by no more than one percentage point and SDR by
less than $0.012$ across their full sweep range. The operating choice
therefore sits inside a flat region for two of the three
parameters, with the GPD threshold the single setting at which a
downstream user would most need to revisit the trade-off.

The sweep script that produced Table~\ref{tab:sensitivity-bali}
accompanies the released code, so the table can be regenerated
without modifying the pipeline.


% ============================================================
\section{Reproducibility Evaluation}\label{sec:results-reproducibility}
% ============================================================

The reproducibility design described in
Section~\ref{sec:methods-reproducibility} is only meaningful if
the pipeline runs end-to-end on commodity hardware. To
demonstrate this, the correction-to-visualisation notebooks were
re-executed on a free-tier Google Colab instance for a
province-scale subdomain of Bali (a $9 \times 14$ pixel grid at
$0.1^{\circ}$, spanning $114.4$ to $115.7^{\circ}$E and $-8.85$
to $-8.05^{\circ}$N), processing all $36$ dekadal windows of the
$2001$ to $2025$ record. Bali was chosen because it is small
enough to run within the free-tier session limits while
exercising every stage of the pipeline.
Appendix~\ref{app:repro}~Figure~\ref{fig:bali-coverage} shows the
subdomain and its data coverage.

Table~\ref{tab:colab-runtime} reports the measured wall-clock
time for the all-dekad batch in each notebook, recorded from the
Colab cell-execution metadata. The correction-to-visualisation
sequence over the
Bali subdomain completed in $72.1$ minutes on a standard Central
Processing Unit (CPU) runtime with no hardware accelerator, well
within a single
free-tier session. The figure covers notebooks $02$ to $06$ and
excludes the two preparatory notebooks, whose timing is dominated
by the domain extent requested and by network throughput and so is
not reproducible.

This $72.1$-minute figure is the cost of all
$36$ dekads for the $126$-pixel Bali grid, not of the full
Indonesian domain, which contains on the order of $10^{4}$ to
$10^{5}$ land pixels and would scale roughly in proportion to the
pixel count for the per-pixel stages. It demonstrates that the
pipeline runs on free-tier infrastructure for a self-contained
region; it is not a timing estimate for a national run. Within
that run, the metrics notebook dominates the budget because it
computes the full $31$-metric catalogue at every pixel and dekad,
and the correction notebook, which trains a per-dekad CNN, is
second. Both were run on the high Random Access Memory (RAM)
runtime Colab offers without charge.

\begin{table}[H]
\centering
\caption[Measured runtime of the LSEQM+DL pipeline on Google Colab]{Measured wall-clock runtime of notebooks $02$ to $06$ of the LSEQM+DL
pipeline on a free-tier Google Colab CPU instance, for the Bali
subdomain across all $36$ dekadal windows\label{tab:colab-runtime}}
\begingroup\small
\begin{tabular}{@{}llc@{}}
\toprule
Notebook & Stage & Batch runtime (min) \\
\midrule
02 & LSEQM+DL bias correction (per-dekad CNN training) & 16.3 \\
03 & Verification metric computation ($31$ metrics)    & 45.4 \\
04 & Quality assessment (CQI sub-scores)               &  6.9 \\
05 & Station validation against BMKG                   &  0.5 \\
06 & Visualisation                                     &  3.0 \\
\midrule
\multicolumn{2}{@{}l}{End-to-end total}                & 72.1 \\
\bottomrule
\end{tabular}
\endgroup
\end{table}

The execution environment was Python $3.12$ on the Colab default
image, with TensorFlow $2.20.0$, NumPy $2.0.2$, and netCDF4
$1.7.4$ among the libraries reported by the run; no pinned
requirements file beyond a single supplementary install was
needed because the remaining dependencies are pre-installed in
the Colab image. The run was a single-pass execution on shared
Colab hardware, so the timings are indicative rather than
benchmarked: the absolute numbers will vary with the runtime
allocation Colab assigns, but the order of magnitude (around an
hour for the Bali subdomain across all dekads, scaling with
pixel count for larger domains) is the operationally relevant
figure.

The present run does not include a formal bit-equivalence test
against the local outputs, so the strongest reproducibility
claim it supports is structural rather than numerical. The
pipeline completes every stage on the free-tier environment and
produces outputs of the expected shape and coverage (for the
Bali grid, $80$ of $126$ land pixels corrected with the
remainder masked as ocean, and $24$ of $24$ contributing CPC
cells fitted). A bit-for-bit comparison cell that asserts
equality of the Colab and local corrected fields within
floating-point tolerance is a small addition that would upgrade
the claim from structural completion to numerical reproduction,
and it is identified here as the next reproducibility deliverable.


% ============================================================
\section*{Chapter summary}
% ============================================================

Against $172$ independent BMKG stations, the LSEQM+DL framework
moves three of the four verification pillars cleanly toward the
gauge target, value adjustment, distribution alignment, and extreme
preservation all reaching near-unity ratios, while the fourth, event
detection, shows a designed trade-off of light-rain hits for a lower
false-alarm ratio. The stage-wise attribution confirms that each
stage moves the metric matched to its design dimension. One
diagnostic stays pinned close to its raw value at every stage: the
daily Pearson correlation holds near $0.35$ per pixel against the
in-sample CPC-UNI target and near $0.28$ per station against the
independent BMKG network, and the same flat pattern holds
month-by-month and at both date labels.

The date label itself,
by contrast, moves the correlation a long way: matching the
satellite to the local-day gauge convention lifts the GPM-era
daily correlation of the raw IMERG-L archive against BMKG, pooled
over all station-days, from $0.20$ to $0.57$. That locates most of
the apparent disagreement in the calendar rather than in the
retrieval, and none of the recovery in the correction, which
leaves $r$ within sampling noise of the raw archive at both
labels. The spatial CQI maps concentrate the skill in the
densely gauged west, with graceful degradation to the LSEQM baseline
in the sparsely gauged east, and the correction-to-visualisation
sequence runs on a free Colab in $72.1$ minutes on the Bali
subdomain.
Chapter~\ref{chap:discussion} explains this flat correlation as a
structural property of marginal correction, not a failure.
